\documentclass[10pt]{article}
\usepackage[T1]{fontenc}
\usepackage[utf8]{inputenc}
\usepackage[french]{babel}
\usepackage{amsmath,amssymb,mathrsfs,mathpazo}
\usepackage{geometry}
\geometry{a4paper,margin=2cm}
\usepackage{graphicx,array,tabularx,tabularray,colortbl,multirow,multicol,booktabs,enumitem,pifont,ulem,fancybox,diagbox,cancel}
\usepackage{pstricks,pstricks-add,pst-plot,pst-tree,pst-node,pst-text,pst-func,pst-math,pst-bspline,pst-3dplot,pst-eucl,pst-solides3d}
\let\pscalculate\undefined
\usepackage{pst-bezier}
\usepackage{tikz,pgfplots,tkz-tab}
\usetikzlibrary{arrows,arrows.meta,patterns,calc,positioning,trees,decorations.markings,shapes.geometric,shapes.misc,shapes.arrows,intersections,decorations.pathmorphing,decorations.pathreplacing,backgrounds,shadows,babel}
\pgfplotsset{compat=1.18}
\usepgfplotslibrary{fillbetween}
\usepackage[np]{numprint}
\usepackage{listings,makecell,pas-tableur,twoopt}
\let\barre\undefined
\usepackage[pstricks]{tablvar}
\definecolor{darkgreen}{RGB}{0,100,0}
\providecommand{\red}{\color{red}}
\providecommand{\blue}{\color{blue}}
\providecommand{\green}{\color{green}}
\providecommand{\black}{\color{black}}
\providecommand{\white}{\color{white}}
\providecommand{\gray}{\color{gray}}
\providecommand{\cyan}{\color{cyan}}
\providecommand{\magenta}{\color{magenta}}
\providecommand{\yellow}{\color{yellow}}
\providecommand{\violet}{\color{violet}}
\providecommand{\orange}{\color{orange}}
\providecommand{\euro}{\texteuro}
\providecommand{\EUR}{\texteuro}
\providecommand{\ieme}{\textsuperscript{e}}
\providecommand{\degres}{\ensuremath{^\circ}}
\providecommand{\degree}{\ensuremath{^\circ}}
\providecommand{\vv}[1]{\overrightarrow{#1}}
\providecommand{\Coord}[1]{#1}
\providecommand{\decosix}{}
\providecommand{\starredbullet}{\ensuremath{\bullet}}
\pagestyle{empty}
\setlength{\parindent}{0pt}
\usepackage{adjustbox}\DeclareUnicodeCharacter{2236}{:}\DeclareUnicodeCharacter{2190}{\ensuremath{\leftarrow}}
\providecommand{\R}{\mathbb{R}}
\providecommand{\N}{\mathbb{N}}
\providecommand{\Z}{\mathbb{Z}}
\providecommand{\Q}{\mathbb{Q}}
\providecommand{\D}{\mathbb{D}}
\providecommand{\C}{\mathbb{C}}
\providecommand{\e}{{\mathrm{e}}}
\providecommand{\vect}[1]{\overrightarrow{#1}}
\providecommand{\vectt}[1]{\overrightarrow{\mathrm{#1}}}
\providecommand{\vvt}[1]{\overrightarrow{\mathrm{#1}}}
\providecommand{\np}[2][]{{#2~#1}}
\providecommand{\vv}[1]{{\overrightarrow{#1}}}
\renewcommand{\d}{{\mathrm{d}}}
\providecommand{\degres}{{^\circ}}
\providecommand{\degre}{{^\circ}}
\providecommand{\blue}{\color{blue}}
\providecommand{\red}{\color{red}}
\providecommand{\enskip}{\quad}
\providecommand{\llbracket}{\lbrack\!\lbrack}
\providecommand{\rrbracket}{\rbrack\!\rbrack}
\providecommand{\euro}{{\text{€}}}
\providecommand{\no}{{\text{n°}}}
\providecommand{\interval}[2]{{[#1;#2]}}
\providecommand{\intervFO}[2]{{[#1;#2[}}
\providecommand{\dotfill}{\ldots}

\providecommand{\ds}{\displaystyle}
\providecommand{\textding}[1]{#1}
\providecommand{\Oij}{\(\mathrm{O};\vec{\imath},\vec{\jmath}\)}
\providecommand{\Oijk}{\(\mathrm{O};\vec{\imath},\vec{\jmath},\vec{k}\)}
\providecommand{\Ouv}{\(\mathrm{O};\vec{u},\vec{v}\)}
\providecommand{\intervFF}[2]{\ensuremath{[#1;#2]}}
\renewcommand{\euro}{\ifmmode\text{\texteuro}\else\texteuro\fi}
\DeclareUnicodeCharacter{00B0}{\ifmmode^{\circ}\else\textdegree\fi}
\DeclareUnicodeCharacter{20AC}{\ifmmode\text{\texteuro}\else\texteuro\fi}
% Énoncés : Éducation nationale. Transcription et corrigés : APMEP, Denis Vergès et contributeurs. Découpage et classement : CoopMaths, Sébastien Lozano, Éric Elter et contributeurs. Adaptation de la présentation : Maths Explorer.
% CC BY-SA : https://coopmaths.fr/www/about/
% https://coopmaths.fr/alea/static/bac/2025/tex/bac_2025_05_sujet1_ameriquenord_4.tex
\begin{document}
\noindent\textit{APMEP / CoopMaths. CC BY-SA. Présentation adaptée par Maths Explorer.}\par\medskip
\medskip

La \textbf{partie C} est indépendante des parties \textbf{A} et \textbf{B}.

\medskip

\textbf{Partie A}

\medskip

On donne ci-dessous, dans un repère orthogonal, les courbes $\mathcal{C}_{1}$ et $\mathcal{C}_{2}$, représentations graphiques de deux fonctions définies et dérivables sur $\R$. L'une des deux fonctions représentées est la fonction dérivée de l'autre. On les notera $g$ et $g’$.

On précise également que :
\begin{itemize}
\item La courbe $\mathcal{C}_{1}$ coupe l'axe des ordonnées au point de coordonnées $(0~;~1)$.
\item La courbe $\mathcal{C}_{2}$ coupe l'axe des ordonnées au point de coordonnées $(0~;~2)$ et l'axe des abscisses aux points de coordonnées $(-2~;~0)$ et $(1~;~0)$.
\end{itemize}

\begin{center}
\begin{tikzpicture}
\begin{axis}[/pgf/number format/.cd, use comma,
x={15mm}, y={6mm}, 
xmin=-3, xmax=5, ymin = -8, ymax= 7, 
xtick = {-2,...,4}, ytick={-7,...,6},
tick label style={font=\footnotesize},
minor tick num = 0, grid=major, axis lines =center]
\node[below left] at (axis cs: 0,0) {{\footnotesize 0}};
\addplot [line width=1.2pt, color=red, smooth, samples=300, domain= -3:5]{(x^2+3*x+1)*exp(-x)};
\node[red] at(axis cs: 2.2,2) {$ \mathcal{C}_1 $};
\addplot [line width=1.2pt, color=blue, dashed, smooth, samples=300, domain= -3:5]{(-x^2-x+2)*exp(-x)};
\node[blue] at(axis cs: 2.2,-1.1) {$ \mathcal{C}_2 $};
\end{axis}
\end{tikzpicture}
















\end{center}

\begin{enumerate}
	\item En justifiant, associer à chacune des fonctions $g$ et $g'$ sa représentation graphique.
	\item Justifier que l'équation réduite de la tangente à la
courbe représentative de la fonction $g$ au point d'abscisse 0 est $y=2x + 1$.
\end{enumerate}

\medskip

\textbf{Partie B}

\medskip

On considère $(E)$ l'équation différentielle 

\[y+y'=(2x + 3) \e^{-x},\]
où $y$ est une fonction de la variable réelle $x$.

\medskip

\begin{enumerate}
\item Montrer que la fonction $f_{0}$ définie pour tout nombre réel $x$ par $f_{0}(x)=\left(x^{2}+3 x\right) \e^{-x}$ est une solution particulière de l'équation différentielle $(E)$.
\item Résoudre l'équation différentielle $\left(E_{0}\right): y+ y'= 0$.
\item Déterminer les solutions de l'équation différentielle $(E)$.
\item On admet que la fonction $g$ décrite dans la \textbf{partie A} est une solution de l'équation différentielle $(E)$.

Déterminer alors l'expression de la fonction $g$.
\item Déterminer les solutions de l'équation différentielle $(E)$ dont la courbe admet exactement deux points d'inflexion.
\end{enumerate}

\medskip

\textbf{Partie C}

\medskip

On considère la fonction $f$ définie pour tout nombre réel $x$ par :

\[f(x)=\left(x^{2}+3 x+2\right) \e^{-x}\]

\begin{enumerate}
\item Démontrer que la limite de la fonction $f$ en $+\infty$ est égale à 0.

On admet par ailleurs que la limite de la fonction $f$ en $-\infty$ est égale à $+\infty$.

\item On admet que la fonction $f$ est dérivable sur $\R$. On note $f’$ la fonction dérivée de $f$ sur $\R$.
	\begin{enumerate}
		\item Vérifier que, pour tout nombre réel $x$, $f’(x)=\left(-x^{2}-x+1\right) \e^{-x}$.
		\item Déterminer le signe de la fonction dérivée $f’$ sur $\R$ puis en déduire les variations de la fonction $f$ sur $\R$.
	\end{enumerate}
\item Expliquer pourquoi la fonction $f$ est positive sur l'intervalle $[0;+\infty[$.
\item On notera $\mathcal{C}_{f}$ la courbe représentative de la fonction $f$ dans un repère orthogonal \Oij. On admet que la fonction $F$ définie pour tout nombre réel $x$ par $F(x)=\left(-x^{2}-5 x-7\right) \e^{-x}$ est une primitive de la fonction $f$.

Soit $\alpha$ un nombre réel positif.

Déterminer l'aire $\mathcal{A}(\alpha)$, exprimée en unité d'aire, du domaine du plan délimité par l'axe des abscisses, la courbe $\mathcal{C}_{f}$ et les droites d'équation $x=0$ et $x=\alpha$.
\end{enumerate}

\end{document}
